\scanpage{209}
\margin{208}
\hfill -- \mspage{51} -- \hfill IX

5. \uline{Le théorème d'orthogonalité dans le cas} $\ell=p$ (\uline{cas semi-stable}). \uline{Dualité des schémas abéliens} $B_o$ \uline{et} $B'_o$. \uline{Caractérisation de la partie fixe}. \uline{Critères de bonne réduction}, \uline{de bonne réduction essentielle et semi-stable}.

5.1. Plaçons nous dans le cas 2.2.3 où $S$ est hensélien et $A^{\unsure{0|o}}$ est $\ell$-divisible, de sorte que la \uline{partie fixe}
\[
T_\ell(A^{\unsure{0|o}})^f\lhook\joinrel\longrightarrow T_\ell(A)
\]
est définie comme le pro-$\ell$-groupe de Barsotti-Tate formé des $({}_{\ell^\nu}A^{\unsure{0|o}})^f$.
Nous allons maintenant donner une définition de la partie torique (2.3.2) de $T_\ell(A^{\unsure{0|o}})$, indépendante de l'hypothèse $\ell\neq p$ faite dans 2.3, mais supposant en revanche $S$ \uline{complet} (*). Pour ceci, introduisons pour tout entier $n\geq0$ le schéma artinien local
\[
\text{(5.1.1)}\qquad S_n=\operatorname{Spec}(V/\underline m^{n+1})\quad ,
\]
n.ème voisinage infinitésimal de $s$ dans $S$, et pour tout $S$-schéma $X$ soit $X_n=X\times_S S_n$. Comme $(A^{\unsure{0|o}})_n=(A_n)^{\unsure{0|o}}$ est lisse sur $S_n$, il est connu (SGA 3 IX 3.6 bis) que $T_o$ est la fibre spéciale d'un unique sous-tore
\[
\text{(5.1.2)}\qquad T_n\subset A_n^{\unsure{0|o}}\quad .
\]
D'après l'unicité des $T_n$, ceux-ci se déduisent les uns des autres par les changements de base $S_{n'}\longrightarrow S_n$. En d'autres termes, il définissent un \uline{sous-tore formel} $\widehat T$ du groupe $\widehat A^{\unsure{0|o}}$ (dans la catégorie des schémas formels sur $S$) complété formel de $A^{\unsure{0|o}}$ le long de la fibre spéciale :
\[
\text{(5.1.3)}\qquad\widehat T\subset\widehat A^{\unsure{0|o}}\quad .
\]
\margin{(*) Cette hypothèse n'est d'ailleurs pas essentielle, comme on montre dans 6.1 ci-dessous.}
\margin{369}
\margin{369}
